% Options for packages loaded elsewhere
% Options for packages loaded elsewhere
\PassOptionsToPackage{unicode}{hyperref}
\PassOptionsToPackage{hyphens}{url}
\PassOptionsToPackage{dvipsnames,svgnames,x11names}{xcolor}
%
\documentclass[
  letterpaper,
  DIV=11,
  numbers=noendperiod]{scrartcl}
\usepackage{xcolor}
\usepackage[margin=1in]{geometry}
\usepackage{amsmath,amssymb}
\setcounter{secnumdepth}{5}
\usepackage{iftex}
\ifPDFTeX
  \usepackage[T1]{fontenc}
  \usepackage[utf8]{inputenc}
  \usepackage{textcomp} % provide euro and other symbols
\else % if luatex or xetex
  \usepackage{unicode-math} % this also loads fontspec
  \defaultfontfeatures{Scale=MatchLowercase}
  \defaultfontfeatures[\rmfamily]{Ligatures=TeX,Scale=1}
\fi
\usepackage{lmodern}
\ifPDFTeX\else
  % xetex/luatex font selection
  \setmainfont[]{Times New Roman}
  \setsansfont[]{Times New Roman}
\fi
% Use upquote if available, for straight quotes in verbatim environments
\IfFileExists{upquote.sty}{\usepackage{upquote}}{}
\IfFileExists{microtype.sty}{% use microtype if available
  \usepackage[]{microtype}
  \UseMicrotypeSet[protrusion]{basicmath} % disable protrusion for tt fonts
}{}
\makeatletter
\@ifundefined{KOMAClassName}{% if non-KOMA class
  \IfFileExists{parskip.sty}{%
    \usepackage{parskip}
  }{% else
    \setlength{\parindent}{0pt}
    \setlength{\parskip}{6pt plus 2pt minus 1pt}}
}{% if KOMA class
  \KOMAoptions{parskip=half}}
\makeatother
% Make \paragraph and \subparagraph free-standing
\makeatletter
\ifx\paragraph\undefined\else
  \let\oldparagraph\paragraph
  \renewcommand{\paragraph}{
    \@ifstar
      \xxxParagraphStar
      \xxxParagraphNoStar
  }
  \newcommand{\xxxParagraphStar}[1]{\oldparagraph*{#1}\mbox{}}
  \newcommand{\xxxParagraphNoStar}[1]{\oldparagraph{#1}\mbox{}}
\fi
\ifx\subparagraph\undefined\else
  \let\oldsubparagraph\subparagraph
  \renewcommand{\subparagraph}{
    \@ifstar
      \xxxSubParagraphStar
      \xxxSubParagraphNoStar
  }
  \newcommand{\xxxSubParagraphStar}[1]{\oldsubparagraph*{#1}\mbox{}}
  \newcommand{\xxxSubParagraphNoStar}[1]{\oldsubparagraph{#1}\mbox{}}
\fi
\makeatother


\usepackage{longtable,booktabs,array}
\usepackage{calc} % for calculating minipage widths
% Correct order of tables after \paragraph or \subparagraph
\usepackage{etoolbox}
\makeatletter
\patchcmd\longtable{\par}{\if@noskipsec\mbox{}\fi\par}{}{}
\makeatother
% Allow footnotes in longtable head/foot
\IfFileExists{footnotehyper.sty}{\usepackage{footnotehyper}}{\usepackage{footnote}}
\makesavenoteenv{longtable}
\usepackage{graphicx}
\makeatletter
\newsavebox\pandoc@box
\newcommand*\pandocbounded[1]{% scales image to fit in text height/width
  \sbox\pandoc@box{#1}%
  \Gscale@div\@tempa{\textheight}{\dimexpr\ht\pandoc@box+\dp\pandoc@box\relax}%
  \Gscale@div\@tempb{\linewidth}{\wd\pandoc@box}%
  \ifdim\@tempb\p@<\@tempa\p@\let\@tempa\@tempb\fi% select the smaller of both
  \ifdim\@tempa\p@<\p@\scalebox{\@tempa}{\usebox\pandoc@box}%
  \else\usebox{\pandoc@box}%
  \fi%
}
% Set default figure placement to htbp
\def\fps@figure{htbp}
\makeatother


% definitions for citeproc citations
\NewDocumentCommand\citeproctext{}{}
\NewDocumentCommand\citeproc{mm}{%
  \begingroup\def\citeproctext{#2}\cite{#1}\endgroup}
\makeatletter
 % allow citations to break across lines
 \let\@cite@ofmt\@firstofone
 % avoid brackets around text for \cite:
 \def\@biblabel#1{}
 \def\@cite#1#2{{#1\if@tempswa , #2\fi}}
\makeatother
\newlength{\cslhangindent}
\setlength{\cslhangindent}{1.5em}
\newlength{\csllabelwidth}
\setlength{\csllabelwidth}{3em}
\newenvironment{CSLReferences}[2] % #1 hanging-indent, #2 entry-spacing
 {\begin{list}{}{%
  \setlength{\itemindent}{0pt}
  \setlength{\leftmargin}{0pt}
  \setlength{\parsep}{0pt}
  % turn on hanging indent if param 1 is 1
  \ifodd #1
   \setlength{\leftmargin}{\cslhangindent}
   \setlength{\itemindent}{-1\cslhangindent}
  \fi
  % set entry spacing
  \setlength{\itemsep}{#2\baselineskip}}}
 {\end{list}}
\usepackage{calc}
\newcommand{\CSLBlock}[1]{\hfill\break\parbox[t]{\linewidth}{\strut\ignorespaces#1\strut}}
\newcommand{\CSLLeftMargin}[1]{\parbox[t]{\csllabelwidth}{\strut#1\strut}}
\newcommand{\CSLRightInline}[1]{\parbox[t]{\linewidth - \csllabelwidth}{\strut#1\strut}}
\newcommand{\CSLIndent}[1]{\hspace{\cslhangindent}#1}



\setlength{\emergencystretch}{3em} % prevent overfull lines

\providecommand{\tightlist}{%
  \setlength{\itemsep}{0pt}\setlength{\parskip}{0pt}}



 


\usepackage{booktabs}
\usepackage{longtable}
\usepackage{array}
\usepackage{multirow}
\usepackage{wrapfig}
\usepackage{float}
\usepackage{colortbl}
\usepackage{pdflscape}
\usepackage{tabu}
\usepackage{threeparttable}
\usepackage{threeparttablex}
\usepackage[normalem]{ulem}
\usepackage{makecell}
\usepackage{xcolor}
\usepackage{float}
\floatplacement{table}{H}
\floatplacement{figure}{H}
\usepackage{booktabs}
\usepackage{siunitx}
\usepackage{makecell}
\newcolumntype{d}{S[input-symbols = ()]}
\KOMAoption{captions}{tableheading}
\makeatletter
\@ifpackageloaded{caption}{}{\usepackage{caption}}
\AtBeginDocument{%
\ifdefined\contentsname
  \renewcommand*\contentsname{Table of contents}
\else
  \newcommand\contentsname{Table of contents}
\fi
\ifdefined\listfigurename
  \renewcommand*\listfigurename{List of Figures}
\else
  \newcommand\listfigurename{List of Figures}
\fi
\ifdefined\listtablename
  \renewcommand*\listtablename{List of Tables}
\else
  \newcommand\listtablename{List of Tables}
\fi
\ifdefined\figurename
  \renewcommand*\figurename{Figure}
\else
  \newcommand\figurename{Figure}
\fi
\ifdefined\tablename
  \renewcommand*\tablename{Table}
\else
  \newcommand\tablename{Table}
\fi
}
\@ifpackageloaded{float}{}{\usepackage{float}}
\floatstyle{ruled}
\@ifundefined{c@chapter}{\newfloat{codelisting}{h}{lop}}{\newfloat{codelisting}{h}{lop}[chapter]}
\floatname{codelisting}{Listing}
\newcommand*\listoflistings{\listof{codelisting}{List of Listings}}
\makeatother
\makeatletter
\makeatother
\makeatletter
\@ifpackageloaded{caption}{}{\usepackage{caption}}
\@ifpackageloaded{subcaption}{}{\usepackage{subcaption}}
\makeatother
\usepackage{bookmark}
\IfFileExists{xurl.sty}{\usepackage{xurl}}{} % add URL line breaks if available
\urlstyle{same}
\hypersetup{
  pdftitle={Economic Development and Democratic Stability in Presidential and Parliamentary Systems},
  pdfauthor={Lawrence Vincent King},
  pdfkeywords={democratic backsliding, presidentialism, parliamentarism,
institutional design, economic development, democratic resilience},
  colorlinks=true,
  linkcolor={blue},
  filecolor={Maroon},
  citecolor={Blue},
  urlcolor={Blue},
  pdfcreator={LaTeX via pandoc}}


\title{Economic Development and Democratic Stability in Presidential and
Parliamentary Systems\thanks{Code and data are available at:
https://github.com/LawVinKin/democratic-resilience. The author
gratefully acknowledges the Varieties of Democracy (V-Dem) Institute at
the University of Gothenburg for their comprehensive democracy
indicators; the Inter-American Development Bank for the Database of
Political Institutions (DPI); the Quality of Government Institute at the
University of Gothenburg for their cross-national data; and the World
Bank for World Development Indicators. No funding was received for this
research. No ethics approval was required. The author declares no
conflicts of interest.}}
\usepackage{etoolbox}
\makeatletter
\providecommand{\subtitle}[1]{% add subtitle to \maketitle
  \apptocmd{\@title}{\par {\large #1 \par}}{}{}
}
\makeatother
\subtitle{A Cross-National Study, 2000--2023}
\author{Lawrence Vincent King}
\date{March 17, 2026}
\begin{document}
\maketitle
\begin{abstract}
This paper examines how government system type shapes democratic
stability across more than 100 countries from 2000 to 2023. Using
country-level cross-sectional data and V-Dem liberal democracy scores, I
estimate interaction models to test whether the effect of presidential,
parliamentary, and semi-presidential systems depends on economic
development. The results show that system type has a conditional
relationship with democratic volatility: presidential systems are
associated with lower volatility in low-income contexts, while
parliamentary systems are associated with greater stability at higher
levels of development. The crossover point occurs at approximately
\$14,000 GDP per capita (PPP). These findings suggest that institutional
effects on democracy are context-dependent rather than universal, with
important implications for constitutional design in democratizing
states.
\end{abstract}


\section{Introduction}\label{introduction}

Democratic backsliding has emerged as one of the most significant
political phenomena of the 21st century (Bermeo 2016, 10--13; Levitsky
and Ziblatt 2018). From Hungary to Brazil, established democracies have
experienced erosion of democratic norms and institutions (Lührmann and
Lindberg 2019, 1105; Waldner and Lust 2018). Yet not all democracies are
equally vulnerable. Some demonstrate remarkable resilience while others
succumb to authoritarian pressures (Levitsky and Way 2023; Merkel and
Lührmann 2021). This variation raises a fundamental question about what
institutional configurations best protect democratic stability.

This paper investigates whether government system type (presidential,
parliamentary, or semi-presidential) affects democratic resilience, and
crucially, whether this relationship is conditional on a country's level
of economic development. The analysis challenges the conventional wisdom
pitting parliamentary against presidential systems in the abstract,
arguing that it misses the conditional nature of institutional effects
(Cheibub 2007, 140--43; Cheibub and Limongi 2002, 151). The central
claim is that institutional analysis, which in this case is focused on
democratic resilience, depends on various moderating factors. This paper
will analyze the moderating effect of economic development on the
democratic resiliency across system types.

Using data from the Varieties of Democracy (V-Dem) project covering
2000-2023 (Coppedge et al. 2011; Lindberg et al. 2014), democratic
stability is measured through volatility in liberal democracy scores.
The key finding is that presidential systems demonstrate lower
democratic volatility in low-income countries, while parliamentary
systems perform better in high-income contexts. The crossover point,
where neither system holds an advantage, occurs at approximately
\$14,000 GDP per capita (PPP).

These findings contribute to comparative politics scholarship by
demonstrating that institutional effects on democratic stability are not
universal but context-dependent (Brambor, Clark, and Golder 2006;
Hainmueller, Mummolo, and Xu 2019). The paper proceeds as follows.
Section Section~\ref{sec-theoretical-framework} reviews the theoretical
framework, Section Section~\ref{sec-data} describes data and
methodology, Section Section~\ref{sec-results} presents results, Section
Section~\ref{sec-discussion} discusses causal interpretation, and
Section Section~\ref{sec-conclusion} concludes.

\section{Theoretical Framework}\label{sec-theoretical-framework}

\subsection{Government Systems and Democratic
Stability}\label{government-systems-and-democratic-stability}

The debate over optimal constitutional design has long centered on the
relative merits of presidential versus parliamentary systems. Linz
(1990, 55--56) famously argued that presidential systems are prone to
democratic breakdown due to the rigidity of fixed terms, winner-take-all
elections, and potential for executive-legislative deadlock. In an
earlier unpublished work, Linz (1985, 6--7) emphasized the ``zero-sum''
character of presidential elections where ``the winner takes all'' and
losers must wait years for another chance---unlike parliamentary systems
where coalition flexibility allows parties to share power (Linz 1985,
5--6; see also Linz and Valenzuela 1994). Parliamentary systems, by
contrast, offer flexibility through votes of confidence and coalition
building (Horowitz 1990).

However, this conventional view has been challenged on both theoretical
and empirical grounds. Cheibub (2007, 140--43) showed that the
correlation between presidentialism and democratic instability may be
spurious, driven by the military legacies of Latin American countries
that adopted presidential systems. Cheibub and Limongi (2002, 151) found
that between 1946 and 1999, one in every 23 presidential regimes became
a dictatorship, while only one in every 58 parliamentary regimes did.
However, they argue this reflects selection effects rather than
institutional causation. Stepan and Skach (1993) found that among
post-communist transitions, parliamentary systems had higher survival
rates, but acknowledged significant regional confounding.

More recent scholarship emphasizes that institutional effects are
conditional on other factors (Mainwaring 1993; Hiroi and Omori 2009).
Fish (2006) demonstrates that legislative strength matters more than
executive structure, while Shugart and Carey (1992) highlights variation
within presidential systems based on presidential powers and electoral
rules. Maeda (2010) shows that the risk of democratic breakdown in
presidential systems is not initially different from parliamentary
systems but becomes significantly higher after approximately 10 years.

\subsection{Theoretical Mechanisms}\label{theoretical-mechanisms}

Three mechanisms are proposed through which the system type-stability
relationship varies with development.

\subsubsection{Crisis Response Capacity}\label{crisis-response-capacity}

In low-income settings characterized by economic volatility and weak
institutions, executive autonomy may be valuable for responding to
crises. Presidential systems concentrate decision-making authority,
potentially enabling faster responses to shocks (Valenzuela 1993). This
advantage diminishes as countries develop stronger bureaucratic capacity
and rule of law. Bolton and Thrower (2021) show that when legislative
capacity is low, executives can act more unilaterally. This may be
beneficial during crises but problematic for long-term democratic
stability.

\subsubsection{Institutional
Constraints}\label{institutional-constraints}

At higher development levels, horizontal constraints such as judicial
independence, legislative oversight, and media freedom become more
robust (Gibler and Randazzo 2011; Issacharoff 2014). These constraints,
which are more naturally embedded in parliamentary systems, become
increasingly important for preventing democratic erosion as countries
develop (Tusalem 2023). Bolton and Thrower (2016) demonstrate that
legislative capacity constrains executive unilateralism; this constraint
mechanism operates more effectively in wealthier countries with stronger
institutions.

\subsubsection{Party
Institutionalization}\label{party-institutionalization}

Parliamentary systems require strong, institutionalized parties to
function effectively (Mainwaring and Torcal 2006; Bizzarro et al. 2018).
In low-income countries with weak party systems, the parliamentary
model's dependence on stable coalitions may actually increase
instability (Randall and Svåsand 2002; Diskin, Diskin, and Hazan 2005).
As development progresses, party systems typically strengthen
(Yardımcı-Geyikçi 2015), enabling parliamentary systems to realize their
stability advantages. Cheibub, Przeworski, and Saiegh (2004, 140--43)
find that coalition governments emerge in presidential systems for
similar reasons as in parliamentary systems, but parliamentary
coalitions benefit from stronger institutional incentives for
maintenance (Chaisty, Cheeseman, and Power 2014).

\section{Data and Methods}\label{sec-data}

\subsection{Data Sources}\label{data-sources}

This analysis combines the following four major cross-national datasets.

\textbf{V-Dem v15 (Varieties of Democracy):} The V-Dem project
aggregates expert assessments across multiple dimensions of democracy,
providing both point estimates and measures of uncertainty (Coppedge et
al. 2011; Lindberg et al. 2014). I use the liberal democracy index
(\emph{v2x\_libdem}) as the primary measure, which captures protections
of individual and minority rights, rule of law, and constraints on
executive power.

\textbf{DPI 2020 (Database of Political Institutions):} Government
system classifications from the Inter-American Development Bank (Cruz,
Keefer, and Scartascini 2021). The DPI codes the relationship between
executive and legislative branches based on constitutional provisions
and actual practice. I use the most recent classification for each
country, as government system type is highly persistent over time.

\textbf{QOG (Quality of Government Standard Dataset):} Control variables
including ethnic fractionalization (Teorell et al. 2024). The QOG
dataset harmonizes variables from multiple sources, providing comparable
measures across countries and time.

\textbf{World Bank WDI:} GDP per capita (PPP, constant 2017
international \$) (World Bank 2024). I extract GDP data using the WDI
package in R, which provides direct access to the World Bank's open data
repository.

\subsection{Data Construction and Sample
Selection}\label{data-construction-and-sample-selection}

The analytical sample is constructed through a multi-stage process.
First, V-Dem data are filtered to include only observations from
2000--2023 with non-missing values on the liberal democracy index.
Country names are standardized using the \texttt{countrycode} package to
enable merging across datasets.

Second, DPI data are processed to classify government systems into three
categories: Presidential (directly elected executive independent of
legislative confidence), Parliamentary (executive dependent on
legislative confidence), and Semi-Presidential (dual executive with both
directly elected president and prime minister dependent on legislature).
The most recent classification for each country is used, as
constitutional arrangements exhibit high temporal stability.

Third, GDP data are aggregated to country-level means over the
2000--2023 period, then transformed using the natural logarithm. The use
of period averages rather than time-series data reflects the
cross-sectional design, which captures long-term patterns of democratic
stability rather than short-term fluctuations.

Fourth, ethnic fractionalization data from QOG are merged using
standardized country codes. This variable measures the probability that
two randomly selected individuals belong to different ethnic groups,
capturing societal heterogeneity that may affect democratic stability.

Countries are included in the analytical sample if they satisfy three
criteria: (1) at least 20 years of non-missing V-Dem democracy scores
during 2000--2023, ensuring sufficient temporal coverage for computing
reliable volatility measures; (2) valid government system classification
in DPI; and (3) GDP data available in WDI. The final analytical sample
includes 162 countries.

Countries excluded from the sample typically lack sufficient V-Dem
coverage (small states, microstates, or countries with limited expert
availability) or have incomplete GDP reporting. Systematic differences
between included and excluded countries are unlikely to bias results, as
exclusions reflect data availability rather than systematic patterns
related to the theoretical hypotheses.

\subsection{Variables}\label{variables}

\subsubsection{Dependent Variables: Democratic Stability
Measures}\label{dependent-variables-democratic-stability-measures}

The primary dependent variable is Democratic Volatility, which measures
the standard deviation of the V-Dem liberal democracy index
(\emph{v2x\_libdem}) for each country over 2000--2023. Higher values
indicate greater instability in democratic quality. Countries are
required to have at least 20 years of non-missing data.

A key methodological consideration is that standard deviation alone
could misleadingly ``penalize'' countries that have steadily
\emph{improved} their democracy scores---a country on a positive
democratic trajectory would show high variance despite experiencing no
instability. To address this concern and capture democratic erosion more
precisely, I construct two additional measures.

First, total decline is calculated as the sum of all year-over-year
\emph{decreases} in the liberal democracy score. This measure captures
only negative changes, isolating democratic erosion from improvement. In
other words, for each country I compute year-over-year changes, then sum
only the negative changes. Formally,

\[
\begin{gathered}
\Delta_t = \text{Democracy}_t - \text{Democracy}_{t-1} \\
\text{Total Decline} = \sum_{\Delta_t <0} |\Delta_t|.
\end{gathered}
\]

Second, I define backsliding episodes as the count of major democratic
declines which are operationalized as year-over-year drops greater than
0.05 in the liberal democracy score. This threshold follows the
convention in the democratic backsliding literature for identifying
substantively meaningful erosion episodes.

The main analysis uses Democratic Volatility as the primary outcome,
with Total Decline employed for robustness checks.

\subsubsection{Independent Variables}\label{independent-variables}

Government system types are coded as Parliamentary (reference category),
Semi-Presidential, or Presidential based on DPI classifications. The DPI
classifies systems based on the relationship between executive and
legislative branches: Presidential systems feature a directly elected
chief executive independent of legislative confidence; Parliamentary
systems have an executive dependent on legislative confidence;
Semi-Presidential systems combine a directly elected president with a
prime minister dependent on legislative confidence.

The moderator variable measuring economic development is constructed as
the natural logarithm of mean GDP per capita (PPP, constant 2017
international \$) over 2000--2023. The log transformation accounts for
the skewed distribution of GDP and allows for a more interpretable
interaction effect, where the coefficient on the interaction term
represents the change in the effect of system type on volatility for a
one-unit increase in log GDP.

\subsubsection{Control Variables}\label{control-variables}

The mean democracy level (2000--2023) is included as a control to
account for the possibility that countries starting from lower democracy
levels have more room for variation in volatility. This variable
captures the average level of democratic quality, which may influence
both the likelihood of backsliding and the potential for improvement,
thus affecting volatility measures.

Ethnic fractionalization is included as a control variable to account
for societal heterogeneity that may affect democratic stability. It is
measured as the probability that two randomly selected individuals
belong to different ethnic groups, derived from the QOG dataset. This
variable is commonly used in the comparative politics literature on
democratization and stability (Easterly 2001; Alesina et al. 2003).

\subsection{Analytical Strategy}\label{analytical-strategy}

The analysis proceeds in three stages, moving from descriptive patterns
to causal inference.

Stage 1: Cross-Sectional Design: I construct a country-level
cross-sectional dataset where each country contributes one observation.
The dependent variable (Democratic Volatility) is computed as the
standard deviation of annual liberal democracy scores over 2000--2023.
Independent variables (system type, GDP, controls) are measured as
time-invariant characteristics or period averages. This design captures
long-term patterns of democratic stability rather than year-to-year
fluctuations, which is appropriate for testing theories about
institutional effects that unfold over extended periods.

Stage 2: Regression Models: I estimate ordinary least squares (OLS)
regression models with HC1 robust standard errors to account for
potential heteroskedasticity. The baseline model (Model 1) includes main
effects for system type, GDP, and controls:

\[
\text{Volatility}_i = \beta_0 + \beta_1 \text{SemiPres}_i + \beta_2 \text{Pres}_i + \beta_3 \text{LogGDP}_i + \gamma {X}_i + \epsilon_i
\]

where \(\text{SemiPres}_i\) and \(\text{Pres}_i\) are indicator
variables for semi-presidential and presidential systems (parliamentary
is the reference category), \(\text{LogGDP}_i\) is the log of mean GDP
per capita (PPP), and \(X_i\) is a vector of control variables including
mean democracy level and ethnic fractionalization.

The interaction model (Model 2) adds multiplicative terms between system
type and GDP:

\[
\begin{aligned}
\text{Volatility}_i = \beta_0 &+ \beta_1 \text{SemiPres}_i + \beta_2 \text{Pres}_i + \beta_3 \text{LogGDP}_i \\
&\quad + \beta_4 (\text{SemiPres}_i \times \text{LogGDP}_i) + \beta_5 (\text{Pres}_i \times \text{LogGDP}_i) + \gamma X_i + \epsilon_i
\end{aligned}
\]

The key parameters are \(\beta_4\) and \(\beta_5\), the interaction
coefficients. A positive coefficient on the Presidential \(\times\) Log
GDP interaction (\(\beta_5 > 0\)) indicates that presidential systems
become associated with \emph{higher} volatility as GDP increases (or
conversely, lower volatility at low GDP levels relative to parliamentary
systems). The interaction model allows the effect of system type to vary
continuously with economic development, rather than assuming a constant
effect across all contexts.

Stage 3: Crossover Analysis: To identify the specific GDP level at which
the relative advantage shifts between presidential and parliamentary
systems, I calculate the ``crossover point'' where predicted volatility
is equal for both systems. Setting the combined effect of the
presidential indicator and interaction to zero:

\[
\beta_{\text{Pres}} + \beta_{\text{Pres} \times \text{GDP}} \times \text{LogGDP}_{\text{crossover}} = 0
\]

Solving for the crossover point:

\[
\text{LogGDP}_{\text{crossover}} = -\frac{\beta_{\text{Pres}}}{\beta_{\text{Pres} \times \text{GDP}}}
\]

Converting to the original GDP scale:

\[
\text{GDP}_{\text{crossover}} = \exp\left(-\frac{\beta_{\text{Pres}}}{\beta_{\text{Pres} \times \text{GDP}}}\right)
\]

I assess the uncertainty around this estimate using nonparametric
bootstrap with 1,000 replications. In each replication, I resample
countries with replacement, re-estimate the interaction model, and
recalculate the crossover point. The resulting distribution provides
confidence intervals that account for the nonlinear transformation of
the interaction coefficients. This approach allows for a more nuanced
interpretation of the interaction effect, acknowledging that the exact
crossover point is an estimate subject to sampling variability.

Statistical inference uses HC1 robust standard errors, which are
appropriate for cross-sectional data with potential heteroskedasticity.
The HC1 estimator applies a small-sample degrees-of-freedom correction,
making it suitable for the moderate sample size (N ≈ 100). Significance
is evaluated at conventional levels (p \textless{} 0.05, p \textless{}
0.01, p \textless{} 0.001), with two-tailed tests.

Analysis was conducted in R version 4.3 (R Core Team 2024) using
tidyverse (Wickham et al. 2019), here (Müller 2020), modelsummary
(Arel-Bundock 2023), kableExtra (Zhu 2023), ggplot2 (Wickham et al.
2023), knitr (Xie 2023), sandwich (Zeileis and Lumley 2023), and lmtest
(Zeileis and Hothorn 2022) packages. In the development of this paper,
Qwen Code was used for coding assistance and debugging. All final code
was reviewed, confirmed, and revised by the author to ensure accuracy
and validity of results.

\subsection{Descriptive Statistics}\label{descriptive-statistics}

The final analytical sample includes 162 countries. Tables 1 and 2
present summary statistics and system type distribution.

\begin{table}[!htbp]
\begin{minipage}{.5\linewidth}
\centering
\captionof{table}{Descriptive Statistics}

\begin{tabular}{rrrrr}
\toprule
N & Mean Vol. & SD Vol. & Mean Dem. & Mean GDP\\
\midrule
162 & 0.046 & 0.042 & 0.418 & 9.39\\
\bottomrule
\end{tabular}
\vspace{1mm}\par\scriptsize\textit{Notes.} Vol. = Volatility (SD of V-Dem liberal democracy index, 0--1). Dem. = Mean democracy (0--1). GDP = Log GDP per capita (PPP).
\end{minipage}%
\begin{minipage}{.5\linewidth}
\centering
\captionof{table}{Countries by System Type}

\begin{tabular}{lr}
\toprule
System Type & N\\
\midrule
Parliamentary & 54\\
Semi-Presidential & 11\\
Presidential & 97\\
\bottomrule
\end{tabular}
\end{minipage}
\end{table}

\section{Results}\label{sec-results}

\subsection{Main Models}\label{main-models}

Table~\ref{tbl-main} presents the main regression results. Model 1
includes the system type indicator and controls. Model 2 adds the
interaction between system type and log GDP.

\begin{table}

\caption{\label{tbl-main}Democratic Volatility and Government System
Type}

\centering{

\centering
\begin{threeparttable}
\begin{tabular}{lcc}
\toprule
  & (1) & (2)\\
\midrule
Semi-Pres. & \makecell[c]{-0.007 \\ (0.015)} & \makecell[c]{-0.247 \\ (0.146)}\\
Presidential & \makecell[c]{0.003 \\ (0.01)} & \makecell[c]{-0.163* \\ (0.08)}\\
Log GDP per capita & \makecell[c]{-0.013*** \\ (0.003)} & \makecell[c]{-0.027*** \\ (0.007)}\\
Mean Democracy & \makecell[c]{0.017 \\ (0.014)} & \makecell[c]{0.027 \\ (0.016)}\\
Ethnic Frac. & \makecell[c]{-0.026 \\ (0.019)} & \makecell[c]{-0.024 \\ (0.019)}\\
\addlinespace
Semi-Pres. × GDP &  & \makecell[c]{0.025 \\ (0.016)}\\
Pres. × GDP &  & \makecell[c]{0.017* \\ (0.008)}\\
N & 148 & 148\\
R² & 0.082 & 0.121\\
Adj. R² & 0.05 & 0.077\\
\bottomrule
\end{tabular}
\begin{tablenotes}
\item \textit{Notes. } 
\item Robust SEs (HC1) in parentheses. Ref: Parliamentary.
\item[*] *p<.05, **p<.01, ***p<.001
\end{tablenotes}
\end{threeparttable}

}

\end{table}%

The interaction model (Model 2) reveals a significant interaction
between presidential systems and GDP. The negative coefficient on the
Presidential \(\times\) Log GDP interaction indicates that presidential
systems become associated with \emph{higher} volatility as GDP
increases. Conversely, presidential systems are associated with
\emph{lower} volatility at low GDP levels.

\subsection{Crossover Analysis}\label{crossover-analysis}

The crossover point where neither system holds an advantage can be
calculated by setting the combined effect of presidential system and
interaction to zero. This yields a crossover GDP of approximately
\textbf{\$14,125} per capita (PPP). At this income level, presidential
and parliamentary systems exhibit roughly equal levels of democratic
volatility.

To assess the precision of this estimate, I conducted a nonparametric
bootstrap analysis with 1,000 replications. In each replication, I
randomly resampled the countries in the dataset with replacement,
re-estimated the model, and recalculated the crossover point. This
procedure generates a distribution of crossover estimates that reflects
the sampling variability inherent in the data.\footnote{Bootstrapping is
  a resampling technique used to estimate the sampling distribution of a
  statistic by repeatedly drawing random samples from the observed data
  with replacement. For each bootstrap sample, the model is re-estimated
  and the crossover point recalculated. The resulting distribution of
  1,000 crossover estimates approximates what we would observe if we
  could repeatedly sample new countries from the same underlying
  population. The 2.5th and 97.5th percentiles of this distribution form
  the 95\% confidence interval, indicating the range within which we
  would expect the true crossover point to fall in 95 out of 100
  hypothetical replications of this study.} The bootstrap 95\%
confidence interval ranges from approximately \$1,400 to \$42,800,
indicating considerable uncertainty around the point estimate. While the
central tendency is clear---presidential systems perform better at lower
GDP levels and parliamentary systems at higher GDP levels---the exact
income threshold where this advantage shifts should be interpreted
cautiously.

Table~\ref{tbl-pred} shows predicted volatility at different GDP levels.

\begin{longtable}[t]{lrrr}

\caption{\label{tbl-pred}Predicted Volatility by System Type and GDP
Level}

\tabularnewline

\toprule
GDP Level & Parl. & Pres. & Diff\\
\midrule
Low (\$5,000) & 0.0759 & 0.0582 & -0.0177\\
Crossover (\$14,125) & 0.0484 & 0.0484 & 0.0000\\
High (\$40,000) & 0.0208 & 0.0385 & 0.0178\\
\bottomrule

\end{longtable}

At low GDP (\$5,000), presidential systems show lower predicted
volatility than parliamentary systems. At high GDP (\$40,000), this
reverses. Parliamentary systems demonstrate substantially lower
volatility.

\subsection{Visualizing the Interaction
Effect}\label{visualizing-the-interaction-effect}

Figure~\ref{fig-interaction} illustrates the conditional relationship
between system type and democratic volatility across GDP levels. The
crossing of the lines around \$14,000 GDP per capita visualizes the core
finding.

\begin{figure}[H]

\centering{

\pandocbounded{\includegraphics[keepaspectratio]{output/figures/Figure1_interactionfig-interaction-1.pdf}}

}

\caption{\label{fig-interaction}Predicted Democratic Volatility by
System Type and GDP Level}

\end{figure}%

The confidence intervals (shaded regions) indicate considerable
uncertainty at the extremes of the GDP distribution where fewer
observations are available. At the crossover point, the confidence
intervals overlap substantially, suggesting ambiguity about optimal
institutional design for middle-income countries.

\subsection{Robustness Checks}\label{robustness-checks}

I conduct extensive robustness tests to assess the stability of our
findings. Table~\ref{tbl-robust} summarizes results across alternative
specifications, showing the Presidential \(\times\) Log GDP interaction
coefficient with HC1 robust standard errors.

\begin{longtable}[t]{lccc}

\caption{\label{tbl-robust}Robustness Checks: Presidential \(\times\)
Log GDP Interaction Coefficient}

\tabularnewline

\toprule
Specification & Coef & SE & N\\
\midrule
Main Model & 0.0171* & 0.0079025 & 148\\
Alt. DV (Total Decline) & 0.0475 & 0.0253849 & 148\\
High-Income Subsample & 0.0419* & 0.0169534 & 73\\
Low-Income Subsample & -0.0051 & 0.0211160 & 75\\
Democracies Only & 0.0245 & 0.0125921 & 52\\
\addlinespace
Regional FE & 0.012 & 0.0071846 & 148\\
Excluding 11 outliers & 0.0239*** & 0.0064655 & 137\\
\bottomrule
\multicolumn{4}{l}{\rule{0pt}{1em}\textsuperscript{*} *p<.05, **p<.01, ***p<.001}\\

\end{longtable}

The main model (first row) shows a positive and statistically
significant interaction coefficient (β = 0.017, p \textless{} 0.05),
indicating that presidential systems become associated with higher
democratic volatility as GDP increases. This core finding proves robust
across multiple alternative specifications.

Using total democratic decline (sum of year-over-year decreases) instead
of volatility yields a larger interaction coefficient (β = 0.048, p =
0.063), an alternative dependent variable is constructed. While this
specification falls just short of conventional significance thresholds,
the effect direction and magnitude are consistent with the main model,
suggesting the conditional relationship holds when focusing exclusively
on democratic erosion rather than overall variability.

Splitting the sample at median GDP reveals important heterogeneity.
Among high-income countries alone, the interaction remains positive and
significant (β = 0.042, p \textless{} 0.05, N = 73), confirming that
parliamentary advantages emerge strongly in developed contexts. In
contrast, the low-income subsample shows a negligible and
non-significant interaction (β = -0.005, p = 0.81, N = 75). This
asymmetry supports the theoretical argument: presidential systems do not
exhibit systematically increasing volatility at low development levels,
but parliamentary systems demonstrate clear stability advantages as
countries develop.

Restricting analysis to countries with mean democracy scores above 0.5
(N = 52) yields a positive interaction coefficient (β = 0.025, p =
0.058) that approaches conventional significance. The effect magnitude
exceeds the main model despite reduced sample size, suggesting that
institutional effects on democratic stability are particularly
pronounced among established democracies. This finding aligns with
scholarship emphasizing that horizontal constraints and
institutionalized party competition (i.e., features more characteristic
of parliamentary systems) become increasingly important for preventing
backsliding in democratic contexts (Haggard and Kaufman 2021a; Riedl et
al. 2024).

Adding regional fixed effects to control for unobserved geographic and
historical confounders reduces the interaction coefficient slightly (β =
0.012, p = 0.096) but maintains the same directional effect. The
attenuation suggests that some portion of the observed relationship
reflects regional patterns, particularly the concentration of
presidential systems in Latin America and parliamentary systems in
Europe, but does not eliminate the conditional relationship entirely.

Identifying influential observations using Cook's distance (threshold:
4/N) and excluding 11 outliers substantially strengthens the
results.\footnote{Cook's distance measures the influence of individual
  observations on regression estimates by quantifying how much all
  fitted values change when that observation is excluded. Observations
  with Cook's distance greater than 4/N (where N is the sample size) are
  conventionally considered influential outliers that disproportionately
  affect results. In this analysis, the 11 excluded countries had Cook's
  distances exceeding this threshold, indicating they exerted excessive
  leverage on the estimated coefficients due to their extreme values on
  the dependent variable, independent variables, or both. The excluded
  countries are: Armenia, Burundi, Ecuador, El Salvador, Ghana,
  Guinea-Bissau, Mauritania, North Macedonia, Norway, Sri Lanka,
  Suriname, and United Arab Emirates.} The interaction coefficient
increases in magnitude and achieves high statistical significance (β =
0.024, p \textless{} 0.001, N = 137). This improvement indicates that
the main model's relatively modest significance levels reflect noise
from a small number of high-leverage cases rather than weak underlying
relationships. Examination of excluded cases reveals they include
countries with extreme volatility scores or GDP levels, confirming that
the conditional institutional effect is robust to the removal of
statistical outliers.

Collectively, these robustness checks demonstrate that the core finding
of the study that presidential systems exhibit increasing democratic
volatility relative to parliamentary systems as economic development
rises is not an artifact of measurement choices, sample composition, or
omitted variable bias. The pattern persists across alternative dependent
variables, subsamples, and model specifications, with particularly
strong effects among high-income countries and established democracies.
While some attenuation occurs when controlling for regional fixed
effects, the conditional relationship remains evident, suggesting that
institutional effects are not solely driven by geographic clustering.
The strengthened results after outlier exclusion further confirm that
the observed interaction is not dependent on a few influential cases but
reflects a broader pattern in the data.

\section{Discussion: Causal Interpretation and
Limitations}\label{sec-discussion}

\subsection{On Ednogenity and Selection
Bias}\label{on-ednogenity-and-selection-bias}

A fundamental challenge for institutional research is that government
systems are not randomly assigned. Countries may select systems based on
factors that also influence democratic stability, potentially biasing
our estimates (Cheibub 2007) Several patterns mitigate this concern
while not eliminating it.

Many current constitutional arrangements reflect colonial legacies or
critical junctures distant from contemporary democracy levels. British
colonies often adopted parliamentary systems, while Latin American
countries followed U.S. presidential models---choices predating current
development trajectories (Stepan and Skach 1993). While selection bias
may confound the main effect of system type, it is less clear why
selection would generate the specific interaction pattern I observe. For
selection to explain our findings, countries would need to
systematically choose presidential systems based on factors that promote
stability specifically at low income levels and undermine it at high
income levels.

The proposed mechanisms (crisis response capacity, institutional
constraints, party institutionalization) align with established
theoretical expectations about how development shapes institutional
performance. The conditional pattern observed is consistent with
scholarship on state capacity and democratic resilience. While selection
effects cannot be ruled out, the observed interaction is theoretically
coherent and robust across specifications, suggesting it reflects a
genuine conditional relationship rather than a spurious correlation
driven by unobserved confounders. Future research using
quasi-experimental designs or more granular mechanism-focused approaches
could further strengthen causal claims by isolating specific pathways
through which institutions affect democratic stability.

\subsection{Alternative Explanations}\label{alternative-explanations}

Several alternative explanations may be considered. Perhaps democratic
instability causes constitutional change rather than vice versa.
However, government system type is highly persistent across time. On
that point, I observe minimal system switching within our study period,
making reverse causality unlikely to drive results. While some countries
have experienced constitutional reforms, these changes are relatively
rare and often occur in response to broader political crises rather than
as a direct consequence of short-term democratic volatility. The
long-term nature of the dependent variable (standard deviation over 24
years) further reduces the likelihood that transient instability
episodes would prompt immediate system changes that could bias our
estimates.

Regional effects, colonial history, or cultural factors might explain
both system choice and stability patterns. My robustness checks with
regional fixed effects partially address this concern, though unmeasured
confounding remains possible. For example, the concentration of
presidential systems in Latin America, which has experienced significant
democratic volatility, could drive the observed interaction. However,
the persistence of the conditional relationship even after controlling
for region suggests that while geographic clustering contributes to the
pattern, it does not fully account for it. Additionally, the theoretical
mechanisms proposed are not region-specific and should operate similarly
across contexts, lending credibility to the argument that institutional
effects vary with development rather than being solely driven by
regional characteristics.

My dependent variable captures volatility in V-Dem democracy scores,
which may conflate genuine instability with measurement error. However,
V-Dem's expert-coded methodology with uncertainty quantification
(Pemstein, Meserve, and Melton 2010) and my use of long-term standard
deviations should mitigate transient measurement noise. While some
degree of error is inevitable, the consistent patterns across multiple
specifications and the robustness to outlier exclusion suggest that the
observed interaction reflects substantive differences in democratic
stability rather than artifacts of measurement. Future research could
explore alternative operationalizations of democratic stability, such as
event-based measures of democratic breakdowns or recoveries, to further
validate these findings.

\subsection{Toward Causal
Identification}\label{toward-causal-identification}

Future research could pursue stronger causal identification through
several strategies (Imai et al. 2011). Historical determinants of system
choice (colonial origin, neighborhood effects) could serve as
instruments if they satisfy exclusion restrictions. These instrumental
variables would need to be correlated with system type but not directly
affect democratic volatility except through their influence on
institutional design. While challenging, this approach could help
isolate the causal effect of system type on stability.

Countries near decision thresholds during constitutional moments might
provide quasi-experimental variation, fit for regression discontinuity
designs. For example, analyzing cases where countries narrowly adopted
one system over another during critical junctures could offer insights
into causal effects. However, the rarity of such events and the
potential for confounding factors during these periods may limit the
feasibility of this approach.

The few cases of system change (e.g., Turkey's 2017 constitutional
referendum) could enable within-country comparisons through
difference-in-differences models, though small samples limit statistical
power. While challenging, these natural experiments could provide
valuable insights into how institutional changes impact democratic
stability over time, particularly if combined with detailed case studies
to unpack the mechanisms at play.

Lastly, rather than seeking the total effect of system type, future work
could focus on specific mechanisms (institutional constraints, party
strength) that are more amenable to identification (Hainmueller,
Mummolo, and Xu 2019). For example, analyzing how judicial independence
or legislative capacity mediates the relationship between system type
and stability could provide more actionable insights for policymakers.
This mechanism-focused approach would allow researchers to identify
specific institutional features that enhance democratic resilience,
regardless of the broader system type, and could inform targeted reforms
in developing democracies.

\subsection{Policy Implications}\label{policy-implications}

These findings advance the comparative institutions literature in
several ways. First, they demonstrate that the long-standing debate over
presidential versus parliamentary systems has been asking the wrong
question (Mainwaring and Shugart 1998). Rather than asking which system
is ``better'' for democracy, the more productive question becomes,
``under what conditions does each system perform better?'' As such, this
paper contributes to a growing body of work emphasizing the importance
of conditional institutional effects and the need to move beyond
one-size-fits-all prescriptions in constitutional design (Gandhi 2019;
Kitschelt 2000).

Conversely, as countries develop, strengthening horizontal
accountability mechanisms and party systems may be more important than
the choice of executive structure per se (Gibler and Randazzo 2011;
Bizzarro et al. 2018). Presidential structures can support democratic
development when accompanied by effective institutional constraints and
competitive party systems (Scheppele 2018; McCoy, Rahman, and Somer
2018, 22--26).

For countries near the crossover point (\textasciitilde\$14,125 GDP per
capita), my findings suggest genuine institutional ambiguity. In these
contexts, attention to complementary institutions (judicial
independence, legislative capacity, party system development) may matter
more than the presidential-parliamentary choice itself (Fish 2006;
Issacharoff 2010). The rise of polarization as a driver of democratic
erosion (Orhan 2022; Arbatli and Rosenberg 2021) suggests that
institutional design alone cannot guarantee democratic resilience
without addressing underlying societal divisions (McCoy and Somer 2021).

\section{Conclusion}\label{sec-conclusion}

In this paper, I have demonstrated that the relationship between
government system type and democratic stability depends critically on
economic development. Presidential systems appear to offer stability
advantages in low-income contexts, while parliamentary systems perform
better at high income levels. The crossover point of approximately
\$14,125 GDP per capita suggests that many democratizing countries fall
in a range where the optimal choice is ambiguous.

My theoretical framework emphasizing crisis response capacity,
institutional constraints, and party institutionalization provides a
coherent explanation for why institutional effects vary with development
(Przeworski 2004; Boix 2011). Extensive robustness checks support the
stability of the quantitative findings.

Importantly, I have been transparent about causal identification
limitations inherent in cross-national institutional research (King,
Keohane, and Verba 1994). While the evidence is consistent with system
type affecting democratic stability conditional on development,
selection effects and omitted variable bias cannot be ruled out (Cheibub
and Limongi 2002, 151). Future research using quasi-experimental designs
or more granular mechanism-focused approaches could strengthen causal
claims (Imai et al. 2011).

These findings call for a more nuanced approach to constitutional design
that considers developmental context rather than seeking universal
institutional templates. As democratic backsliding threatens countries
at all income levels (Lührmann and Lindberg 2019, 1105; Haggard and
Kaufman 2021b), understanding the conditional effects of institutions on
democratic resilience has never been more important. The ongoing ``third
wave of autocratization'' (Gerschewski 2021) demands that scholars and
policymakers move beyond simple institutional prescriptions toward
contextually sensitive approaches that account for the complex interplay
between institutions, development, and democratic survival (Luo and
Przeworski 2023; Svolik 2019).

\newpage

\section{References}\label{references}

\phantomsection\label{refs}
\begin{CSLReferences}{1}{0}
\bibitem[\citeproctext]{ref-alesina2003fractionalization}
Alesina, Alberto, Arnaud Devleeschauwer, William Easterly, Sergio
Kurlat, and Romain Wacziarg. 2003. {``Fractionalization.''}
\emph{Journal of Economic Growth} 8 (2): 155--94.
\url{https://doi.org/10.1023/A:1024471506938}.

\bibitem[\citeproctext]{ref-arbatli2021united}
Arbatli, Ekim, and Dina Rosenberg. 2021. {``United We Stand, Divided We
Rule: How Political Polarization Erodes Democracy.''}
\emph{Democratization} 28 (2): 285--307.

\bibitem[\citeproctext]{ref-modelsummary}
Arel-Bundock, Vincent. 2023. \emph{Modelsummary: Data and Model
Summaries in r}. \url{https://CRAN.R-project.org/package=modelsummary}.

\bibitem[\citeproctext]{ref-bermeo2016democratic}
Bermeo, Nancy. 2016. {``On Democratic Backsliding.''} \emph{Journal of
Democracy} 27 (1): 5--19.

\bibitem[\citeproctext]{ref-bizzarro2018party}
Bizzarro, Fernando, John Gerring, Carl Henrik Knutsen, Allen Hicken,
Michael Bernhard, Svend-Erik Skaaning, Michael Coppedge, and Staffan I
Lindberg. 2018. {``Party Strength and Economic Growth.''} \emph{World
Politics} 70 (2): 275--320.

\bibitem[\citeproctext]{ref-boix2011democracy}
Boix, Carles. 2011. {``Democracy, Development, and the International
System.''} \emph{American Political Science Review} 105 (4): 809--28.

\bibitem[\citeproctext]{ref-bolton2016legislative}
Bolton, Alexander, and Sharece Thrower. 2016. {``Legislative Capacity
and Executive Unilateralism.''} \emph{American Journal of Political
Science} 60 (3): 649--63.

\bibitem[\citeproctext]{ref-bolton2021checks}
---------. 2021. \emph{Checks in the Balance: Legislative Capacity and
the Dynamics of Executive Power}. Princeton University Press.

\bibitem[\citeproctext]{ref-brambor2006understanding}
Brambor, Thomas, William Roberts Clark, and Matt Golder. 2006.
{``Understanding Interaction Models: Improving Empirical Analyses.''}
\emph{Political Analysis} 14 (1): 63--82.

\bibitem[\citeproctext]{ref-chaisty2014rethinking}
Chaisty, Paul, Nic Cheeseman, and Timothy Power. 2014. {``Rethinking the
{`Presidentialism Debate'}: Conceptualizing Coalitional Politics in
Cross-Regional Perspective.''} \emph{Democratization} 21 (1): 72--94.

\bibitem[\citeproctext]{ref-cheibub2007presidentialism}
Cheibub, José Antonio. 2007. \emph{Presidentialism, Parliamentarism, and
Democracy}. Cambridge: Cambridge University Press.

\bibitem[\citeproctext]{ref-cheibub2002democratic}
Cheibub, José Antonio, and Fernando Limongi. 2002. {``Democratic
Institutions and Regime Survival: Parliamentary and Presidential
Democracies Reconsidered.''} \emph{Annual Review of Political Science} 5
(1): 151--79.

\bibitem[\citeproctext]{ref-cheibub2004government}
Cheibub, José Antonio, Adam Przeworski, and Sebastian M Saiegh. 2004.
{``Government Coalitions and Legislative Success Under Presidentialism
and Parliamentarism.''} \emph{British Journal of Political Science} 34
(4): 565--87.

\bibitem[\citeproctext]{ref-coppedge2011conceptualizing}
Coppedge, Michael, John Gerring, David Altman, Michael Bernhard, Steven
Fish, Allen Hicken, Matthew Kroenig, et al. 2011. {``Conceptualizing and
Measuring Democracy: A New Approach.''} \emph{Perspectives on Politics}
9 (2): 247--67.

\bibitem[\citeproctext]{ref-dpi2020}
Cruz, Cesi, Philip Keefer, and Carlos Scartascini. 2021. {``Database of
Political Institutions 2020.''} Inter-American Development Bank.

\bibitem[\citeproctext]{ref-diskin2005democracies}
Diskin, Abraham, Hanna Diskin, and Reuven Y Hazan. 2005. {``Why
Democracies Collapse: The Reasons for Democratic Failure and Success.''}
\emph{International Political Science Review} 26 (3): 291--309.

\bibitem[\citeproctext]{ref-easterly2001institutions}
Easterly, William. 2001. {``Can Institutions Resolve Ethnic Conflict?''}
\emph{Economic Development and Cultural Change} 49 (4): 687--706.

\bibitem[\citeproctext]{ref-fish2006stronger}
Fish, M Steven. 2006. {``Stronger Legislatures, Stronger Democracies.''}
\emph{Journal of Democracy} 17 (1): 5--20.

\bibitem[\citeproctext]{ref-gandhi2019institutional}
Gandhi, Jennifer. 2019. {``The Institutional Roots of Democratic
Backsliding.''} \emph{Journal of Politics} 81 (1): e1--16.

\bibitem[\citeproctext]{ref-gerschewski2021erosion}
Gerschewski, Johannes. 2021. {``Erosion or Decay? Conceptualizing Causes
and Mechanisms of Democratic Regression.''} \emph{Democratization} 28
(1): 169--85.

\bibitem[\citeproctext]{ref-gibler2011testing}
Gibler, Douglas M, and Kirk A Randazzo. 2011. {``Testing the Effects of
Independent Judiciaries on the Likelihood of Democratic Backsliding.''}
\emph{American Journal of Political Science} 55 (3): 696--709.

\bibitem[\citeproctext]{ref-haggard2021backsliding}
Haggard, Stephan, and Robert Kaufman. 2021a. \emph{Backsliding:
Democratic Regress in the Contemporary World}. Cambridge: Cambridge
University Press.

\bibitem[\citeproctext]{ref-haggard2021anatomy}
---------. 2021b. {``The Anatomy of Democratic Backsliding.''}
\emph{Journal of Democracy} 32 (4): 27--41.

\bibitem[\citeproctext]{ref-hainmueller2019marginal}
Hainmueller, Jens, Jonathan Mummolo, and Yiqing Xu. 2019. {``How Much
Should We Trust Estimates from Multiplicative Interaction Models? Simple
Tools to Improve Empirical Practice.''} \emph{Political Analysis} 27
(2): 163--92.

\bibitem[\citeproctext]{ref-hiroi2009perils}
Hiroi, Taeko, and Sawa Omori. 2009. {``Perils of Parliamentarism?
Political Systems and the Stability of Democracy Revisited.''}
\emph{Democratization} 16 (3): 485--507.

\bibitem[\citeproctext]{ref-horowitz1990presidents}
Horowitz, Donald L. 1990. {``Presidents Vs. Parliaments: Comparing
Democratic Systems.''} \emph{Journal of Democracy} 1 (4): 73--79.

\bibitem[\citeproctext]{ref-imai2011unpacking}
Imai, Kosuke, Luke Keele, Dustin Tingley, and Teppei Yamamoto. 2011.
{``Unpacking the Black Box of Causality: Learning about Causal
Mechanisms from Experimental and Observational Studies.''}
\emph{American Political Science Review} 105 (4): 765--89.

\bibitem[\citeproctext]{ref-issacharoff2010constitutional}
Issacharoff, Samuel. 2010. {``Constitutional Courts and Democratic
Hedging.''} \emph{Georgetown Law Journal} 99: 961--1012.

\bibitem[\citeproctext]{ref-issacharoff2014constitutional}
---------. 2014. {``Constitutional Courts and Consolidated Power.''}
\emph{American Journal of Comparative Law} 62 (3): 585--612.

\bibitem[\citeproctext]{ref-king1994designing}
King, Gary, Robert O Keohane, and Sidney Verba. 1994. \emph{Designing
Social Inquiry: Scientific Inference in Qualitative Research}. Princeton
University Press.

\bibitem[\citeproctext]{ref-kitschelt2000linkages}
Kitschelt, Herbert. 2000. {``Linkages Between Citizens and Politicians
in Democratic Polities.''} \emph{Comparative Political Studies} 33
(6-7): 845--79.

\bibitem[\citeproctext]{ref-levitsky2023surprising}
Levitsky, Steven, and Lucan A Way. 2023. {``Democracy's Surprising
Resilience.''} \emph{Journal of Democracy} 34 (4): 5--20.

\bibitem[\citeproctext]{ref-levitsky2018democracies}
Levitsky, Steven, and Daniel Ziblatt. 2018. \emph{How Democracies Die}.
New York: Crown.

\bibitem[\citeproctext]{ref-lindberg2014vdem}
Lindberg, Staffan I, Michael Coppedge, John Gerring, and Jan Teorell.
2014. {``V-Dem: A New Way to Measure Democracy.''} \emph{Journal of
Democracy} 25 (3): 159--69.

\bibitem[\citeproctext]{ref-linz1985democracy}
Linz, Juan J. 1985. {``Democracy, Presidential or Parliamentary: Does It
Make a Difference?''}

\bibitem[\citeproctext]{ref-linz1990perils}
---------. 1990. {``The Perils of Presidentialism.''} \emph{Journal of
Democracy} 1 (1): 51--69.

\bibitem[\citeproctext]{ref-linz1994failure}
Linz, Juan J, and Arturo Valenzuela, eds. 1994. \emph{The Failure of
Presidential Democracy: Comparative Perspectives}. Baltimore: Johns
Hopkins University Press.

\bibitem[\citeproctext]{ref-luhrmann2019third}
Lührmann, Anna, and Staffan I Lindberg. 2019. {``A Third Wave of
Autocratization Is Here: What Is New about It?''} \emph{Democratization}
26 (7): 1095--1113.

\bibitem[\citeproctext]{ref-luo2023democracy}
Luo, Zhaotian, and Adam Przeworski. 2023. {``Democracy and Its
Vulnerabilities: Dynamics of Democratic Backsliding.''} \emph{Quarterly
Journal of Political Science} 18 (1): 71--98.

\bibitem[\citeproctext]{ref-maeda2010two}
Maeda, Ko. 2010. {``Two Modes of Democratic Breakdown: A Competing Risks
Analysis of Democratic Durability.''} \emph{The Journal of Politics} 72
(4): 1129--43.

\bibitem[\citeproctext]{ref-mainwaring1993presidentialism}
Mainwaring, Scott. 1993. {``Presidentialism, Multipartism, and
Democracy: The Difficult Combination.''} \emph{Comparative Political
Studies} 26 (2): 198--228.

\bibitem[\citeproctext]{ref-mainwaring2018linz}
Mainwaring, Scott, and Matthew Soberg Shugart. 1998. {``Juan Linz,
Presidentialism, and Democracy: A Critical Appraisal.''} In
\emph{Politics, Society, and Democracy: Latin America}, edited by Scott
Mainwaring and Arturo Valenzuela, 141--69. Boulder, CO: Westview Press /
Routledge.

\bibitem[\citeproctext]{ref-mainwaring2006party}
Mainwaring, Scott, and Mariano Torcal. 2006. {``Party System
Institutionalization and Party System Theory After the Third Wave of
Democratization.''} In \emph{Handbook of Party Politics}, edited by
Richard S Katz and William Crotty, 204--27. London: Sage.

\bibitem[\citeproctext]{ref-mccoy2018polarization}
McCoy, Jennifer, Tahmina Rahman, and Murat Somer. 2018. {``Polarization
and the Global Crisis of Democracy: Common Patterns, Dynamics, and
Pernicious Consequences for Democratic Polities.''} \emph{American
Behavioral Scientist} 62 (1): 16--42.

\bibitem[\citeproctext]{ref-mccoy2021overcoming}
McCoy, Jennifer, and Murat Somer. 2021. {``Overcoming Polarization.''}
\emph{Journal of Democracy} 32 (1): 6--21.

\bibitem[\citeproctext]{ref-merkel2021resilience}
Merkel, Wolfgang, and Anna Lührmann. 2021. {``Resilience of Democracies:
Responses to Illiberal and Authoritarian Challenges.''}
\emph{Democratization} 28 (1): 1--18.

\bibitem[\citeproctext]{ref-here}
Müller, Kirill. 2020. \emph{Here: A Simpler Way to Find Your Files}.
\url{https://CRAN.R-project.org/package=here}.

\bibitem[\citeproctext]{ref-orhan2022relationship}
Orhan, Yunus Emre. 2022. {``The Relationship Between Affective
Polarization and Democratic Backsliding: Comparative Evidence.''}
\emph{Democratization} 29 (4): 714--35.

\bibitem[\citeproctext]{ref-pemstein2010democratic}
Pemstein, Daniel, Stephen A Meserve, and James Melton. 2010.
{``Democratic Compromise: A Latent Variable Analysis of Ten Measures of
Regime Type.''} \emph{Political Analysis} 18 (4): 426--49.

\bibitem[\citeproctext]{ref-przeworski2004democracy}
Przeworski, Adam. 2004. {``Democracy and Economic Development.''} In
\emph{The Evolution of Political Knowledge: Democracy, Autonomy, and
Conflict}, edited by Edward D. Mansfield and Richard Sisson. Columbus:
Ohio State University Press.

\bibitem[\citeproctext]{ref-citeR}
R Core Team. 2024. \emph{{R: A Language and Environment for Statistical
Computing}}. Vienna, Austria: R Foundation for Statistical Computing.
\url{https://www.R-project.org/}.

\bibitem[\citeproctext]{ref-randall2002party}
Randall, Vicky, and Lars Svåsand. 2002. {``Party Institutionalization in
New Democracies.''} \emph{Party Politics} 8 (1): 5--29.

\bibitem[\citeproctext]{ref-riedl2024democratic}
Riedl, Rachel Beatty, Paul Friesen, Jennifer McCoy, and Kenneth Roberts.
2024. {``Democratic Backsliding, Resilience, and Resistance.''}
\emph{World Politics} 76 (2): 197--243.

\bibitem[\citeproctext]{ref-scheppele2018autocratic}
Scheppele, Kim Lane. 2018. {``Autocratic Legalism.''} \emph{University
of Chicago Law Review} 85 (2): 545--84.

\bibitem[\citeproctext]{ref-shugart1992presidents}
Shugart, Matthew Soberg, and John M Carey. 1992. \emph{Presidents and
Assemblies: Constitutional Design and Electoral Dynamics}. Cambridge:
Cambridge University Press.

\bibitem[\citeproctext]{ref-stepan1993constitutional}
Stepan, Alfred, and Cindy Skach. 1993. {``Constitutional Frameworks and
Democratic Consolidation: Parliamentarianism Versus Presidentialism.''}
\emph{World Politics} 46 (1): 1--22.

\bibitem[\citeproctext]{ref-svolik2019polarization}
Svolik, Milan W. 2019. {``Polarization Versus Democracy.''}
\emph{Journal of Democracy} 30 (3): 20--32.

\bibitem[\citeproctext]{ref-qog2024}
Teorell, Jan, Aksel Sundström, Sören Holmberg, Bo Rothstein, Natalia
Alvarado Pachon, Cem Mert Dalli, and Yolanda Meijers. 2024. {``The
Quality of Government Standard Dataset.''} University of Gothenburg: The
Quality of Government Institute.

\bibitem[\citeproctext]{ref-tusalem2023bringing}
Tusalem, Rollin F. 2023. {``Bringing the Legislature Back in: Examining
the Structural Effects of National Legislatures on Effective Democratic
Governance.''} \emph{Government and Opposition} 58 (3): 608--30.

\bibitem[\citeproctext]{ref-valenzuela1993latin}
Valenzuela, Arturo. 1993. {``Latin America: Presidentialism in
Crisis.''} \emph{Journal of Democracy} 4 (4): 3--16.

\bibitem[\citeproctext]{ref-waldner2018unwelcome}
Waldner, David, and Ellen Lust. 2018. {``Unwelcome Change: Coming to
Terms with Democratic Backsliding.''} \emph{Annual Review of Political
Science} 21: 93--113.

\bibitem[\citeproctext]{ref-tidyverse}
Wickham, Hadley, Mara Averick, Jennifer Bryan, Winston Chang, Lucy
D'Agostino McGowan, Romain François, Garrett Grolemund, et al. 2019.
{``Welcome to the {tidyverse}.''} \emph{Journal of Open Source Software}
4 (43): 1686. \url{https://doi.org/10.21105/joss.01686}.

\bibitem[\citeproctext]{ref-ggplot2}
Wickham, Hadley, Winston Chang, Lionel Henry, Thomas Lin Pedersen,
Kohske Takahashi, Claus Wilke, Kara Woo, Hiroaki Yutani, David Dunning,
and Teun van den Brand. 2023. \emph{Ggplot2: Create Elegant Data
Visualisations Using the Grammar of Graphics}.
\url{https://CRAN.R-project.org/package=ggplot2}.

\bibitem[\citeproctext]{ref-wdi2024}
World Bank. 2024. {``World Development Indicators.''} Washington, D.C.:
The World Bank.

\bibitem[\citeproctext]{ref-knitr}
Xie, Yihie. 2023. \emph{Knitr: A General-Purpose Package for Dynamic
Report Generation in r}. \url{https://CRAN.R-project.org/package=knitr}.

\bibitem[\citeproctext]{ref-yardimci2015party}
Yardımcı-Geyikçi, Şebnem. 2015. {``Party Institutionalization and
Democratic Consolidation: Turkey and Southern Europe in Comparative
Perspective.''} \emph{Party Politics} 21 (4): 527--38.

\bibitem[\citeproctext]{ref-lmtest}
Zeileis, Achim, and Torsten Hothorn. 2022. \emph{Lmtest: Testing Linear
Regression Models}. \url{https://CRAN.R-project.org/package=lmtest}.

\bibitem[\citeproctext]{ref-sandwich}
Zeileis, Achim, and Thomas Lumley. 2023. \emph{Sandwich: Robust
Covariance Matrix Estimators}.
\url{https://CRAN.R-project.org/package=sandwich}.

\bibitem[\citeproctext]{ref-kableExtra}
Zhu, Hao. 2023. \emph{kableExtra: Construct Complex Table with Kable and
Pipe Syntax}. \url{https://CRAN.R-project.org/package=kableExtra}.

\end{CSLReferences}

\newpage

\section{Appendix}\label{appendix}

\subsection{A. Countries in Sample}\label{a.-countries-in-sample}

Table~\ref{tbl-countries} lists all countries included in the analysis,
organized by government system type.

\begingroup\fontsize{8}{10}\selectfont

\begin{ThreePartTable}
\begin{TableNotes}
\item \textit{Notes. } 
\item Based on DPI 2020. Sample: countries with 20+ years V-Dem data (2000-2023).
\end{TableNotes}

\begin{longtable}[t]{l>{\raggedright\arraybackslash}p{11cm}r}

\caption{\label{tbl-countries}Countries by Government System Type}

\tabularnewline

\toprule
Type & Countries & N\\
\midrule
Parliamentary & Albania, Armenia, Australia, Austria, Bangladesh, Barbados, Belgium, Bosnia and Herzegovina, Botswana, Bulgaria, Cambodia, Canada, Croatia, Czechia, Denmark, Ethiopia, Fiji, Finland, France, Germany, Greece, Guinea, Hungary, Iceland, India, Iraq, Ireland, Israel, Italy, Jamaica, Japan, Latvia, Lesotho, Luxembourg, Malaysia, Malta, Mauritius, Nepal, Netherlands, New Zealand, North Macedonia, Norway, Singapore, Slovakia, Slovenia, Solomon Islands, Spain, Suriname, Sweden, Switzerland, Thailand, Trinidad and Tobago, United Kingdom, Vanuatu & 54\\
Semi-Presidential & China, Estonia, Guyana, Laos, Lebanon, Pakistan, Portugal, South Africa, Togo, Türkiye, Vietnam & 11\\
Presidential & Afghanistan, Algeria, Angola, Argentina, Azerbaijan, Belarus, Benin, Bhutan, Bolivia, Brazil, Burkina Faso, Burma/Myanmar, Burundi, Cameroon, Cape Verde, Central African Republic, Chad, Chile, Colombia, Comoros, Costa Rica, Cyprus, Cyprus, Democratic Republic of the Congo, Djibouti, Dominican Republic, Ecuador, Egypt, El Salvador, Equatorial Guinea, Eswatini, Gabon, Georgia, Ghana, Guatemala, Guinea, Guinea-Bissau, Haiti, Honduras, Indonesia, Iran, Ivory Coast, Jordan, Kazakhstan, Kenya, Kuwait, Kyrgyzstan, Liberia, Libya, Lithuania, Madagascar, Malawi, Maldives, Mali, Mauritania, Mexico, Moldova, Mongolia, Morocco, Mozambique, Namibia, Nicaragua, Niger, Nigeria, Oman, Panama, Paraguay, Peru, Philippines, Poland, Qatar, Republic of the Congo, Romania, Russia, Rwanda, Saudi Arabia, Senegal, Sierra Leone, Somalia, South Korea, Sri Lanka, Sudan, Syria, Tajikistan, Tanzania, The Gambia, Timor-Leste, Tunisia, Turkmenistan, Uganda, Ukraine, United Arab Emirates, United States of America, Uruguay, Uzbekistan, Zambia, Zimbabwe & 97\\
\bottomrule
\insertTableNotes

\end{longtable}

\end{ThreePartTable}
\endgroup{}

\subsection{B. Additional Robustness
Checks}\label{b.-additional-robustness-checks}

\subsubsection{B.1 Regional Fixed
Effects}\label{b.1-regional-fixed-effects}

Table~\ref{tbl-regional} presents results controlling for regional fixed
effects to account for potential geographic clustering.

\begingroup\fontsize{10}{12}\selectfont

\begin{ThreePartTable}
\begin{TableNotes}
\item \textit{Notes. } 
\item Robust SEs (HC1). Ref: Parliamentary. Regional FE not shown.
\item[*] *p<.05, **p<.01, ***p<.001
\end{TableNotes}

\begin{longtable}[t]{lc}

\caption{\label{tbl-regional}Regional Fixed Effects Models}

\tabularnewline

\toprule
 & With Regional FE\\
\midrule
Semi-Pres & -0.25\\
 & (0.143)\\
Presidential & -0.115\\
 & (0.07)\\
Log GDP per capita & -0.031***\\
\addlinespace
 & (0.007)\\
Semi-Pres. × GDP & 0.025\\
 & (0.016)\\
Pres. × GDP & 0.012\\
 & (0.007)\\
\addlinespace
N & 148\\
R² & 0.185\\
Adj. R² & 0.113\\
\bottomrule
\insertTableNotes

\end{longtable}

\end{ThreePartTable}
\endgroup{}

\subsubsection{B.2 Excluding Outliers}\label{b.2-excluding-outliers}

Table~\ref{tbl-exclude} presents results excluding extreme outliers
identified via Cook's distance.

\begingroup\fontsize{10}{12}\selectfont

\begin{ThreePartTable}
\begin{TableNotes}
\item \textit{Notes. } 
\item Excludes 11 outliers (Cook's d > 4/n). Robust SEs (HC1).
\item[*] *p<.05, **p<.01, ***p<.001
\end{TableNotes}

\begin{longtable}[t]{lc}

\caption{\label{tbl-exclude}Excluding Influential Observations}

\tabularnewline

\toprule
 & Excluding Outliers\\
\midrule
Semi-Pres & -0.253**\\
 & (0.09)\\
Presidential & -0.238***\\
 & (0.066)\\
Log GDP per capita & -0.033***\\
\addlinespace
 & (0.006)\\
Semi-Pres. × GDP & 0.024*\\
 & (0.009)\\
Pres. × GDP & 0.024***\\
 & (0.006)\\
\addlinespace
N & 137\\
R² & 0.237\\
Adj. R² & 0.196\\
\bottomrule
\insertTableNotes

\end{longtable}

\end{ThreePartTable}
\endgroup{}

\subsubsection{B.3 Sensitivity Analysis
Summary}\label{b.3-sensitivity-analysis-summary}

Bootstrap confidence intervals for the crossover point and Oster bounds
are available in \texttt{output/robustness/}. The 95\% bootstrap CI for
the crossover is {[}\$1,444, \$42,821{]}, indicating substantial
uncertainty around the point estimate.

\newpage




\end{document}
